% chapter-01.tex
\chapter{Introduction}

\section{Background and Motivation}

Biomedical knowledge production has reached a scale that challenges conventional approaches to evidence access and synthesis. Clinicians, biomedical researchers, and applied data scientists routinely work in domains where relevant findings are distributed across large and constantly expanding publication streams. Under these conditions, the practical bottleneck is no longer simple document availability. The dominant challenge is the ability to identify relevant evidence, map it to a specific question, and preserve traceability between retrieved sources and generated analytical conclusions.

Traditional keyword search systems remain valuable for broad exploration and bibliographic retrieval, but they tend to degrade when user intent is semantically rich, context dependent, or expressed in non-standard terminology. In biomedical settings, this limitation is amplified by terminological variation, evolving nomenclature, and heterogeneous reporting styles across journals and subdisciplines. A question phrased from a clinical perspective may require retrieval of evidence described with mechanistic terminology, and vice versa. As a result, lexical matching alone often produces either incomplete recall or noisy result sets that increase downstream interpretation cost.

At the same time, standalone generative methods have introduced a second class of limitations. A model can produce fluent and coherent responses while remaining insufficiently anchored to verifiable evidence. For general-purpose tasks this may be acceptable in exploratory settings; for biomedical decision support it is not. Biomedical workflows require justifiable outputs, explicit provenance, and reproducible behavior under controlled conditions. Consequently, retrieval quality and generation quality cannot be evaluated independently if the system is intended for evidence-grounded use.

Retrieval Augmented Generation (RAG) has emerged as a response to this gap by coupling retrieval and synthesis within a single inference pipeline. In this paradigm, the model response is conditioned on context retrieved at query time, rather than relying exclusively on parametric memory. This coupling creates a stronger basis for verifiability because candidate evidence can be inspected and linked to answer content. For biomedical informatics, this is a meaningful shift: it aligns computational response generation with core requirements of evidence-aware reasoning, source transparency, and update sensitivity.

The motivation of this thesis is therefore not limited to building another search interface or another generative application. The central motivation is to define and evaluate a system architecture that can sustain continuous biomedical ingestion, semantically robust retrieval, and explicit evaluation of retrieval-generation behavior under reproducible experimental conditions. This motivation places architecture, operations, and evaluation methodology on equal footing. In this view, model quality is necessary but insufficient without lifecycle control and traceable measurement.

\section{Problem Statement}

This thesis addresses the integrated problem of continuous knowledge acquisition, structured indexing, and reproducible evaluation for biomedical question answering supported by retrieval augmentation. The emphasis on integration is essential. Each stage can be implemented in isolation, yet isolated optimization frequently produces systems that are difficult to validate, hard to compare, and fragile under operational change.

The first dimension of the problem concerns knowledge freshness. Biomedical evidence changes continuously, and any static corpus progressively diverges from current literature. A platform that cannot update systematically will eventually produce answers based on stale context, even when retrieval and generation components are individually strong. Freshness, therefore, is not a convenience feature; it is a validity requirement.

The second dimension concerns representational suitability for retrieval. New documents must be transformed into indexable units that preserve enough semantic signal for meaningful nearest-neighbor retrieval while remaining computationally tractable. This requires decisions about document granularity, metadata retention, and update semantics. Inadequate representation introduces hidden bias into retrieval, which then propagates into synthesis quality.

The third dimension concerns empirical validity. A system that produces compelling demonstrations but lacks reproducible evaluation procedures cannot support rigorous comparative analysis. Retrieval and generation behavior must be measured through repeatable experimental stages with preserved intermediate artifacts, so that observed performance can be interpreted as a property of the system configuration rather than an artifact of uncontrolled execution.

Existing approaches commonly privilege one dimension at the expense of the others. Static retrieval systems tend to support stable indexing but offer limited mechanisms for continuous acquisition. Standalone generative systems can deliver fluent responses but weaken evidence accountability when retrieval grounding is absent or inconsistent. Ad hoc evaluation scripts may provide snapshots of performance but often fail to preserve sufficient process traceability for robust replication. The central research problem of this thesis is therefore architectural: how to design a biomedical platform that unifies continuous ingestion, retrieval-grounded synthesis, and reproducible evaluation without collapsing these requirements into disconnected tooling.

\section{Research Objectives}

The first objective is to design a modular biomedical RAG architecture that separates acquisition, retrieval, synthesis, and evaluation concerns while preserving explicit integration boundaries. Modularity is treated as a methodological objective rather than only a software engineering preference. It enables controlled replacement of components, supports clearer attribution of observed behavior, and reduces confounding effects during comparative experimentation.

The second objective is to enable automated literature ingestion from PubMed through scheduled execution. This objective includes systematic handling of document identity through Digital Object Identifier (DOI) and PubMed Identifier (PMID) workflows, because identifier consistency directly affects deduplication, update tracking, and downstream reproducibility. Automated ingestion is framed here as a mechanism for maintaining evidence continuity, not merely as a convenience for operational automation.

The third objective is to define a reproducible evaluation framework for both retrieval and generation behavior. This objective requires stage-level artifact production, explicit run organization, and metrics-oriented execution paths that allow reanalysis. The intent is to ensure that reported behavior can be interpreted through transparent process history rather than isolated endpoint measurements.

The fourth objective is to support evidence-grounded question answering with configurable retrieval strategies, including Hypothetical Document Embeddings (HyDE), while preserving observability of system behavior. This objective introduces a deliberate tradeoff: advanced retrieval augmentation can improve semantic recall in difficult queries, yet it may increase latency variance and dependence on additional generation steps. The thesis therefore treats configurability and measurement as inseparable.

Collectively, these objectives prioritize scientific validity, reproducibility, and architectural clarity over isolated model-centric optimization. The expected outcome is a thesis-level research artifact in which design decisions can be connected to measurable system behavior with defensible methodological grounding.

\section{Contributions of the Thesis}

This thesis makes four principal contributions within the scope defined above. First, it contributes a system architecture for biomedical RAG that treats ingestion, retrieval, synthesis, and evaluation as coordinated subsystems rather than disconnected modules. The contribution is architectural: it specifies how component boundaries and interaction contracts can support both operational execution and research-grade analysis.

Second, it contributes an automated scheduling approach for recurrent PubMed ingestion. This contribution operationalizes continuous acquisition and directly addresses corpus staleness as a structural risk in biomedical question answering. By integrating scheduling with downstream ingestion workflows, the thesis frames freshness management as an architectural responsibility.

Third, it contributes a reproducible evaluation framework that structures retrieval and generation analysis into staged experimental pipelines. This contribution is methodological: it enables performance interpretation through preserved intermediate outputs and repeatable execution logic, thereby strengthening the internal validity of empirical conclusions.

Fourth, it contributes an empirical analysis framework for evidence-grounded answer generation under realistic operational constraints. The contribution is not presented as a claim of universal superiority over alternative systems; rather, it provides a controlled basis for discussing how architectural choices influence retrieval behavior, synthesis quality, and reproducibility characteristics.

These contributions are deliberately framed with conservative scope. The thesis does not claim to eliminate all limitations of biomedical question answering, nor does it claim generalization beyond the evaluated setup. Its contribution is a rigorously structured foundation for subsequent extension, replication, and comparative study.

\section{Thesis Structure}

The remainder of the thesis is organized to move from conceptual grounding to architectural definition, then to implementation and empirical analysis. Chapter 2 reviews related literature and establishes the scientific context for retrieval-augmented biomedical question answering. It clarifies how existing retrieval and generation paradigms motivate the architectural direction adopted in this work.

Chapter 3 presents the system architecture, including subsystem responsibilities, workflow design, and operational constraints. It explains why the platform is structured as coordinated services and how architectural tradeoffs are handled with respect to reproducibility, observability, and evidence grounding.

Chapter 4 describes implementation-level realization of the architecture. This chapter maps conceptual components to concrete service behavior and runtime orchestration choices, with emphasis on how the implemented system preserves the design intentions introduced earlier.

Chapter 5 defines the experimental methodology. It specifies evaluation objectives, experimental stages, and measurement strategy, so that empirical claims in later chapters can be interpreted within a transparent protocol rather than as isolated outcomes.

Chapter 6 reports evaluation results and examines retrieval and generation behavior through the defined methodology. The chapter emphasizes interpretation of system behavior in relation to architectural and pipeline decisions.

Chapter 7 provides discussion, including implications for biomedical information access and system design. It connects empirical findings to broader methodological questions about evidence grounding, reproducibility, and operational tradeoffs.

Chapter 8 addresses limitations and future work. It distinguishes current system boundaries from prospective extensions and clarifies which constraints arise from architecture, external dependencies, or evaluation scope.

Chapter 9 concludes the thesis by summarizing the problem framing, contributions, and principal findings. It restates the significance of integrating continuous ingestion, retrieval-grounded synthesis, and reproducible evaluation in a single biomedical platform.

This chapter progression is designed to preserve argumentative continuity: each chapter establishes prerequisites for the next, ensuring that final conclusions are grounded in explicit architectural reasoning and empirically traceable methodology.
